% mazzi: 19  % del 19 solo le pp. 7--12
\chapter{La capacità di carico}
\label{cap:capacitadicarico}

\textbf{Capacità di carico}: caratteristica dell'organismo grazie alla quale i carichi fisici e
psichici che esso è in grado di realizzare possono essere \textbf{rielaborati} --- cioè possono dar
luogo a quello che in metodologia dell'allenamento si definisce \textbf{adattamento} --- senza che
si producano alterazioni della \textbf{salute} e dell'\textbf{allenabilità}.

L'esposizione segue \textbf{Fronher G.} È il passaggio dal dato della prestazione ai concetti con
cui impostare l'allenamento: la match analysis del cap.~\ref{cap:matchanalysis} misura
\textit{quanto} carico la gara impone, la capacità di carico riguarda \textit{quanto} carico
l'organismo può rielaborare. Non va dimenticato che \textbf{la partita è il massimo carico} negli
sport di squadra --- non solo dal punto di vista cinematico e fisiologico, ma anche
\textbf{psicologico e di stress}.

\begin{center}
\begin{forest} classalbero
[Capacità\\di carico
  [Capacità generale di carico\\dell'organismo]
  [Capacità di carico\\meccanico]
  [Capacità di carico dei sistemi\\determinanti la prestazione]
  [Capacità di carico\\psicosociale]
]
\end{forest}
\end{center}

Delle quattro, quella dei \textbf{sistemi determinanti la prestazione} è la ragione per cui questo
argomento arriva qui e non altrove: si può comprendere \textbf{solo} avendo capito quali siano le
dinamiche cinematiche e fisiologiche della prestazione, che è esattamente il lavoro fatto nel
capitolo precedente.

% ---------------------------------------------------------------------------
\section{La capacità generale di carico dell'organismo}

\rubrica{Che cos'è}
Corrisponde essenzialmente allo \textbf{stato generale di salute e di allenabilità}. È espressione
della rielaborazione del carico da parte dell'\textbf{intero organismo} e comprende la
\textbf{capacità di recupero} dopo carichi di qualità e quantità diverse.

\rubrica{Da che cosa è caratterizzata}
\begin{itemize}
  \item dallo stato dei \textbf{sistemi globali di regolazione}:
    \begin{itemize}
      \item sistema nervoso vegetativo;
      \item sistema ormonale;
      \item sistema neuromuscolare;
    \end{itemize}
  \item dalla \textbf{funzionalità dei processi fisiologici e psichici di base}.
\end{itemize}

\rubrica{Come si manifestano le alterazioni}
Quando si va oltre la capacità generale di carico:
\begin{itemize}
  \item aumento della \textbf{frequenza delle malattie infettive};
  \item \textbf{alterazione degli elettroliti};
  \item \textbf{stanchezza generale} --- da non confondere con l'\textit{overtraining}, che è un
    quadro diverso e non trattato qui.
\end{itemize}

% ---------------------------------------------------------------------------
\section{La capacità di carico meccanico}

\rubrica{Che cos'è}
Ne fanno parte quelle \textbf{condizioni biologiche} che, soddisfacendo ciò che è necessario per
\textbf{tollerare carichi meccanici}, ne permettono la \textbf{ripetizione senza problemi}.

È la componente \textbf{più semplice da gestire} ed è propriamente il terreno del preparatore, che
sottopone l'atleta a un \textbf{carico esterno} di natura meccanica --- sia il lavoro in campo, sia
il lavoro con sovraccarico in palestra.

\rubrica{Da che cosa dipende}
Dalla \textbf{sintonia strutturale e funzionale} tra:
\begin{itemize}
  \item il \textbf{sistema neuromuscolare};
  \item le \textbf{componenti passive} dell'apparato locomotorio e di sostegno.
\end{itemize}

\rubrica{Come si manifestano le alterazioni}
\begin{itemize}
  \item \textbf{dolori alle inserzioni tendinee};
  \item \textbf{alterazioni posturali};
  \item \textbf{squilibri muscolari} con le relative sindromi dolorose,
\end{itemize}
che possono poi sfociare in \textbf{traumi e infortuni}.

% ---------------------------------------------------------------------------
\section{La capacità di carico dei sistemi che determinano la prestazione}

\rubrica{Che cos'è}
I \textbf{sistemi e le funzioni principali} necessari per ottenere una \textbf{prestazione sportiva
specifica}.

\rubrica{Perché è centrale}
\begin{itemize}
  \item deve essere il \textbf{parametro di riferimento} per il processo di allenamento sportivo
    specifico. \textbf{Senza conoscere le sollecitazioni di gara non si può programmare
    l'allenamento}, né scegliere i mezzi più funzionali al miglioramento della performance --- sul
    piano cinematico come su quello fisiologico;
  \item per questa ragione, nello \textbf{sport di alto livello}, i programmi di allenamento sono
    prevalentemente diretti a queste condizioni biologiche e sportive;
  \item i suoi adattamenti si raggiungono con \textbf{maggior sicurezza} rispetto a quelli della
    capacità generale di carico e della capacità di carico meccanico.
\end{itemize}

\rubrica{Come si manifestano le alterazioni}
Nei carichi di resistenza si esprimono soprattutto in cambiamenti:
\begin{itemize}
  \item del \textbf{metabolismo};
  \item della \textbf{regolazione cardiocircolatoria};
  \item della \textbf{capacità funzionale dei muscoli}.
\end{itemize}

\rubrica{Un esempio: la mobilità articolare}
Il giocatore di calcio a 5 necessita di una grande \textbf{mobilità articolare}
(cap.~\ref{cap:introfutsal}), richiesta dai cambi di direzione e dalle situazioni con palla ---
finte e cambi di appoggio podalico (cap.~\ref{cap:tecnicatattica}). Andare oltre la capacità di
carico di questi sistemi può \textbf{compromettere le capacità funzionali} dei muscoli interessati.

% ---------------------------------------------------------------------------
\section{La capacità di carico psicosociale}

È la componente \textbf{meno misurabile} di tutte, e la slide non ne dà una definizione: presenta un
punto interrogativo circondato dai fattori che vi concorrono, lasciando l'elenco esplicitamente
aperto. Riguarda tutti gli aspetti \textbf{non direttamente relativi all'allenamento}, di cui però
il preparatore e lo staff dovrebbero tenere conto.

\begin{itemize}
  \item gli \textbf{aumenti improvvisi del carico} di allenamento, che possono provocare al
    giocatore uno stress non necessario;
  \item il \textbf{regime quotidiano di vita} e il \textbf{sonno};
  \item le \textbf{dinamiche familiari}: la disponibilità all'allenamento di un giocatore può
    dipendere da fattori esterni al campo e alla squadra;
  \item le \textbf{richieste esterne troppo elevate}, cioè aspettative altrui eccessive --- meno
    nello sport per adulti, molto di più nello sport giovanile;
  \item un \textbf{regime di vita da ``atleta''} eccessivamente rigido.
\end{itemize}

% ---------------------------------------------------------------------------
\section{Restare dentro la capacità di carico è prevenzione}

È la ragione per cui l'argomento viene trattato:

\begin{quote}
la \textbf{prima azione di prevenzione reale} che il preparatore atletico, come l'allenatore, può
svolgere è \textbf{rimanere all'interno della capacità di carico del soggetto}.
\end{quote}

Restando dentro quella capacità e non andando mai oltre, si sta già svolgendo un'azione di
prevenzione degli infortuni e dei traumi a cui il giocatore potrebbe andare incontro.

\rubrica{Che cosa occorre conoscere}
Per esercitare invece un'\textbf{azione positiva} sull'organismo attraverso lo sport:
\begin{enumerate}
  \item le \textbf{funzioni e le caratteristiche biologiche} dei tessuti di cui sono costituiti i
    sistemi che si vanno a sollecitare --- nel nostro caso non si può prescindere dalla
    \textbf{biologia e fisiologia muscolare} (cap.~\ref{cap:neuromuscolare});
  \item gli \textbf{stimoli specifici} che possono far sviluppare un determinato sistema;
  \item le \textbf{reazioni specifiche} di quei sistemi a quegli stimoli;
  \item i \textbf{segnali e i sintomi} di un carico sbagliato o eccessivo, così da limitare i
    danni;
  \item la \textbf{dinamica del recupero}.
\end{enumerate}
Conoscendo la richiesta di gara si può agire in maniera specifica sul \textbf{singolo individuo}, e
migliorarne così la performance.
